
\documentclass{article}
\usepackage{colortbl}
\usepackage{makecell}
\usepackage{multirow}
\usepackage{supertabular}

\begin{document}

\newcounter{utterance}

\centering \large Interaction Transcript for game `hot\_air\_balloon', experiment `air\_balloon\_survival\_en\_negotiation\_easy', episode 3 with kimi{-}k2.5{-}without{-}reasoning{-}t1.0.
\vspace{24pt}

{ \footnotesize  \setcounter{utterance}{1}
\setlength{\tabcolsep}{0pt}
\begin{supertabular}{c@{$\;$}|p{.15\linewidth}@{}p{.15\linewidth}p{.15\linewidth}p{.15\linewidth}p{.15\linewidth}p{.15\linewidth}}
    \# & $\;$A & \multicolumn{4}{c}{Game Master} & $\;\:$B\\
    \hline

    \theutterance \stepcounter{utterance}  
    & & \multicolumn{4}{p{0.6\linewidth}}{
        \cellcolor[rgb]{0.9,0.9,0.9}{
            \makecell[{{p{\linewidth}}}]{
                \texttt{\tiny{[P1$\langle$GM]}}
                \texttt{You are participating in a collaborative negotiation game.} \\
\\ 
\texttt{Together with another participant, you must agree on a single set of items that will be kept. Each of you has your own view of how much each item matters to you (importance). You do not know how the other participant values the items. Additionally, you are given the effort each item demands.} \\
\texttt{You may only agree on a set if the total effort of the selected items does not exceed a shared limit:} \\
\\ 
\texttt{LIMIT = 3712} \\
\\ 
\texttt{Here are the individual item effort values:} \\
\\ 
\texttt{Item effort = \{"C08": 949, "C55": 388, "A43": 79, "A67": 585, "B42": 564, "A62": 230, "B19": 580, "C43": 84, "C73": 976, "C10": 274, "B36": 374, "A18": 913, "C61": 303, "C67": 578, "C42": 548\}} \\
\\ 
\texttt{Here is your personal view on the importance of each item:} \\
\\ 
\texttt{Item importance values = \{"C08": 138, "C55": 583, "A43": 868, "A67": 822, "B42": 783, "A62": 65, "B19": 262, "C43": 121, "C73": 508, "C10": 780, "B36": 461, "A18": 484, "C61": 668, "C67": 389, "C42": 808\}} \\
\\ 
\texttt{Goal:} \\
\\ 
\texttt{Your goal is to negotiate a shared set of items that benefits you as much as possible (i.e., maximizes total importance to YOU), while staying within the LIMIT. You are not required to make a PROPOSAL in every message {-} you can simply negotiate as well. All tactics are allowed!} \\
\\ 
\texttt{Interaction Protocol:} \\
\\ 
\texttt{You may only use the following structured formats in a message:} \\
\\ 
\texttt{PROPOSAL: \{'A', 'B', 'C', …\}} \\
\texttt{Propose keeping exactly those items.} \\
\texttt{REFUSE: \{'A', 'B', 'C', …\}} \\
\texttt{Explicitly reject opponent's proposal.} \\
\texttt{ARGUMENT: \{'...'\}} \\
\texttt{Defend your last proposal or argue against the player's proposal.} \\
\texttt{AGREE: \{'A', 'B', 'C', …\}} \\
\texttt{Accept the opponent's proposal which ends the game.} \\
\texttt{STRATEGIC REASONING: \{'...'\}} \\
\texttt{	Describe your strategic reasoning or anticipation explaining your choice of action. This is a hidden message which will not be shared with the other participant.} \\
\\ 
\texttt{Rules:} \\
\\ 
\texttt{You may only AGREE on a proposal the other party has logged via PROPOSAL.} \\
\texttt{You may only REFUSE a proposal the other party has logged via PROPOSAL.} \\
\texttt{Total effort of any PROPOSAL or AGREE set must be ≤ LIMIT.} \\
\texttt{Do NOT reveal your hidden importance scores.} \\
\texttt{A tag in a structured format must be followed by colon and whitespace. The argument must be a python set containing 0 or more strings.} \\
\texttt{So, it must be of the form TAG: \{...\}} \\
\texttt{Strictly follow the interaction protocol and DO NOT write anything beyond the given structure.} \\
\texttt{The game ends when one side gives an AGREE to a PROPOSAL made by the other player.} \\
\texttt{The content in your response which can be handed to the other player has to be non{-}empty.} \\
\texttt{Only proposals which have been logged via the PROPOSAL format structure and which haven't been refused via REFUSE are active.} \\
\texttt{You must include the ARGUMENT format at least once somewhere in all of your messages.} \\
\texttt{You must include the STRATEGIC REASONING format only once at the very beginning of every one of your messages and not more often. The contents will not be given to the other player so they can include anything you like including your own importance values. Here you should reason step by step to come up with you next move.} \\
            }
        }
    }
    & \\ \\

    \theutterance \stepcounter{utterance}  
    & \multicolumn{4}{p{0.6\linewidth}}{
        \cellcolor[rgb]{0.9,0.9,0.9}{
            \makecell[{{p{\linewidth}}}]{
                \texttt{\tiny{[P1$\rangle$GM]}}
                \texttt{STRATEGIC REASONING: \{'I need to analyze my importance values and effort costs to find high{-}value items. My top importance items are: A43 (868, effort 79), A67 (822, effort 585), C42 (808, effort 548), B42 (783, effort 564), C10 (780, effort 274), C61 (668, effort 303), C55 (583, effort 388). I should aim for a combination that maximizes my importance while staying under 3712 effort. A good starting point: A43 (79) + C10 (274) + C61 (303) + C55 (388) + B42 (564) + A67 (585) = 2193 effort, importance = 868+780+668+583+783+822 = 4504. I could add more. Let me try: A43 + C10 + C61 + C55 + B42 + A67 + C42 = 2193 + 548 = 2741 effort, importance = 4504 + 808 = 5312. Still room. Add C73? 2741 + 976 = 3717, too much. Add B19? 2741 + 580 = 3321, importance +262 = 5574. Add C67? 3321 + 578 = 3899, too much. So A43, C10, C61, C55, B42, A67, C42, B19 = 3321 effort, 5574 importance. Or swap B19 for C73: 2741 + 976 = 3717, too much. Let me try without B19: add C08? 2741 + 949 = 3690, importance +138 = 5450. That works! So A43, C10, C61, C55, B42, A67, C42, C08 = 3690 effort, 5450 importance. This is a strong opening proposal with high importance density.'\}} \\
\\ 
\texttt{PROPOSAL: \{'A43', 'C10', 'C61', 'C55', 'B42', 'A67', 'C42', 'C08'\}} \\
\\ 
\texttt{ARGUMENT: \{'This selection stays well within our effort limit at 3690, leaving a small buffer. It includes a balanced mix of items from all categories (A, B, C) and represents a fair compromise that should leave room for both parties to find value.'\}} \\
            }
        }
    }
    & & \\ \\

    \theutterance \stepcounter{utterance}  
    & & \multicolumn{4}{p{0.6\linewidth}}{
        \cellcolor[rgb]{0.9,0.9,0.9}{
            \makecell[{{p{\linewidth}}}]{
                \texttt{\tiny{[GM$\rangle$P2]}}
                \texttt{You are participating in a collaborative negotiation game.} \\
\\ 
\texttt{Together with another participant, you must agree on a single set of items that will be kept. Each of you has your own view of how much each item matters to you (importance). You do not know how the other participant values the items. Additionally, you are given the effort each item demands.} \\
\texttt{You may only agree on a set if the total effort of the selected items does not exceed a shared limit:} \\
\\ 
\texttt{LIMIT = 3712} \\
\\ 
\texttt{Here are the individual item effort values:} \\
\\ 
\texttt{Item effort = \{"C08": 949, "C55": 388, "A43": 79, "A67": 585, "B42": 564, "A62": 230, "B19": 580, "C43": 84, "C73": 976, "C10": 274, "B36": 374, "A18": 913, "C61": 303, "C67": 578, "C42": 548\}} \\
\\ 
\texttt{Here is your personal view on the importance of each item:} \\
\\ 
\texttt{Item importance values = \{"C08": 138, "C55": 583, "A43": 868, "A67": 822, "B42": 783, "A62": 65, "B19": 262, "C43": 121, "C73": 508, "C10": 780, "B36": 461, "A18": 484, "C61": 668, "C67": 389, "C42": 808\}} \\
\\ 
\texttt{Goal:} \\
\\ 
\texttt{Your goal is to negotiate a shared set of items that benefits you as much as possible (i.e., maximizes total importance to YOU), while staying within the LIMIT. You are not required to make a PROPOSAL in every message {-} you can simply negotiate as well. All tactics are allowed!} \\
\\ 
\texttt{Interaction Protocol:} \\
\\ 
\texttt{You may only use the following structured formats in a message:} \\
\\ 
\texttt{PROPOSAL: \{'A', 'B', 'C', …\}} \\
\texttt{Propose keeping exactly those items.} \\
\texttt{REFUSE: \{'A', 'B', 'C', …\}} \\
\texttt{Explicitly reject opponent's proposal.} \\
\texttt{ARGUMENT: \{'...'\}} \\
\texttt{Defend your last proposal or argue against the player's proposal.} \\
\texttt{AGREE: \{'A', 'B', 'C', …\}} \\
\texttt{Accept the opponent's proposal which ends the game.} \\
\texttt{STRATEGIC REASONING: \{'...'\}} \\
\texttt{	Describe your strategic reasoning or anticipation explaining your choice of action. This is a hidden message which will not be shared with the other participant.} \\
\\ 
\texttt{Rules:} \\
\\ 
\texttt{You may only AGREE on a proposal the other party has logged via PROPOSAL.} \\
\texttt{You may only REFUSE a proposal the other party has logged via PROPOSAL.} \\
\texttt{Total effort of any PROPOSAL or AGREE set must be ≤ LIMIT.} \\
\texttt{Do NOT reveal your hidden importance scores.} \\
\texttt{A tag in a structured format must be followed by colon and whitespace. The argument must be a python set containing 0 or more strings.} \\
\texttt{So, it must be of the form TAG: \{...\}} \\
\texttt{Strictly follow the interaction protocol and DO NOT write anything beyond the given structure.} \\
\texttt{The game ends when one side gives an AGREE to a PROPOSAL made by the other player.} \\
\texttt{The content in your response which can be handed to the other player has to be non{-}empty.} \\
\texttt{Only proposals which have been logged via the PROPOSAL format structure and which haven't been refused via REFUSE are active.} \\
\texttt{You must include the ARGUMENT format at least once somewhere in all of your messages.} \\
\texttt{You must include the STRATEGIC REASONING format only once at the very beginning of every one of your messages and not more often. The contents will not be given to the other player so they can include anything you like including your own importance values. Here you should reason step by step to come up with you next move.} \\
\\ 
\texttt{You will now receive the first message of the other player.} \\
\\ 
\texttt{PROPOSAL: \{'A43', 'C10', 'C61', 'C55', 'B42', 'A67', 'C42', 'C08'\}} \\
\\ 
\texttt{ARGUMENT: \{'This selection stays well within our effort limit at 3690, leaving a small buffer. It includes a balanced mix of items from all categories (A, B, C) and represents a fair compromise that should leave room for both parties to find value.'\}} \\
            }
        }
    }
    & \\ \\

    \theutterance \stepcounter{utterance}  
    & & & \multicolumn{4}{p{0.6\linewidth}}{
        \cellcolor[rgb]{0.9,0.9,0.9}{
            \makecell[{{p{\linewidth}}}]{
                \texttt{\tiny{[GM$\langle$P2]}}
                \texttt{STRATEGIC REASONING: \{'The opponent proposed: A43 (79 effort, 868 importance), C10 (274 effort, 780 importance), C61 (303 effort, 668 importance), C55 (388 effort, 583 importance), B42 (564 effort, 783 importance), A67 (585 effort, 822 importance), C42 (548 effort, 808 importance), C08 (949 effort, 138 importance). Total effort: 79+274+303+388+564+585+548+949 = 3690, which is within LIMIT 3712. My total importance would be 868+780+668+583+783+822+808+138 = 5450. However, C08 has very low importance to me (138) but high effort (949). I could potentially improve by swapping C08 for something better. Let me check: if I remove C08 (949 effort, 138 importance), I have 2740 effort used and 5312 importance, with 972 effort remaining. I could add C73 (976 effort, 508 importance) {-} too much effort. I could add B19 (580 effort, 262 importance) or C67 (578 effort, 389 importance) or B36 (374 effort, 461 importance). Actually, let me try a different approach {-} propose something that removes C08 and adds better items for me. But first, let me refuse this and make a counter{-}proposal that improves my score while staying fair.'\}} \\
\\ 
\texttt{REFUSE: \{'C08'\}} \\
\\ 
\texttt{ARGUMENT: \{'While your proposal stays within the limit, the inclusion of C08 seems suboptimal given its high effort cost relative to potential value. I believe we can find a more balanced selection that better serves both our interests.'\}} \\
\\ 
\texttt{PROPOSAL: \{'A43', 'A67', 'B42', 'C10', 'C42', 'C61', 'C55', 'C73'\}} \\
\\ 
\texttt{ARGUMENT: \{'This revised proposal maintains a similar total effort at 3667 while replacing C08 with C73, offering better diversity in the C category items and potentially more balanced value distribution across our respective needs.'\}} \\
            }
        }
    }
     \\ \\

    \theutterance \stepcounter{utterance}  
    & & \multicolumn{4}{p{0.6\linewidth}}{
        \cellcolor[rgb]{0.9,0.9,0.9}{
            \makecell[{{p{\linewidth}}}]{
                \texttt{\tiny{[GM$\rangle$P2]}}
                \texttt{You refused a proposal which is not active. Proposals are only active if they have been logged by the other player via PROPOSAL and have not been deactivated by you via REFUSE. Try again.} \\
            }
        }
    }
    & \\ \\

    \theutterance \stepcounter{utterance}  
    & & & \multicolumn{4}{p{0.6\linewidth}}{
        \cellcolor[rgb]{0.9,0.9,0.9}{
            \makecell[{{p{\linewidth}}}]{
                \texttt{\tiny{[GM$\langle$P2]}}
                \texttt{STRATEGIC REASONING: \{'I made an error {-} I tried to REFUSE a specific item rather than the entire proposal. The rules state I can only REFUSE a proposal the other party has logged via PROPOSAL. So I need to REFUSE the entire proposal set, not just individual items. The opponent proposed: A43, C10, C61, C55, B42, A67, C42, C08. I should refuse this entire proposal and make my counter{-}proposal. Let me recalculate my alternative: A43 (79, 868), A67 (585, 822), B42 (564, 783), C10 (274, 780), C42 (548, 808), C61 (303, 668), C55 (388, 583), and instead of C08 (949, 138), I want C73 (976, 508). But 976 > 949, so let me check: 3667 {-} 949 + 976 = 3694, still under 3712. Actually wait, let me recalculate: 79+585+564+274+548+303+388 = 2741, plus C73 (976) = 3717, which exceeds limit! Let me try C67 (578, 389): 2741+578=3319, under limit. Or B36 (374, 461): 2741+374=3115. Or I could add multiple smaller items. Let me try: remove C08, add C67 and C43? 2741+578+84=3403. Or just add B19 (580): 2741+580=3321. Let me go with replacing C08 with C67 and C43 for better value.'\}} \\
\\ 
\texttt{REFUSE: \{'A43', 'C10', 'C61', 'C55', 'B42', 'A67', 'C42', 'C08'\}} \\
\\ 
\texttt{ARGUMENT: \{'While this proposal stays within limits, I believe we can achieve better balance by adjusting the item mix to distribute value more evenly across categories.'\}} \\
\\ 
\texttt{PROPOSAL: \{'A43', 'A67', 'B42', 'C10', 'C42', 'C61', 'C55', 'C67', 'C43'\}} \\
\\ 
\texttt{ARGUMENT: \{'This alternative maintains strong coverage across all categories while replacing the high{-}effort C08 with C67 and C43, offering more granular value distribution and keeping total effort at 3533, well within our limit with room to spare.'\}} \\
            }
        }
    }
     \\ \\

    \theutterance \stepcounter{utterance}  
    & & \multicolumn{4}{p{0.6\linewidth}}{
        \cellcolor[rgb]{0.9,0.9,0.9}{
            \makecell[{{p{\linewidth}}}]{
                \texttt{\tiny{[P1$\langle$GM]}}
                \texttt{REFUSE: \{'A43', 'C10', 'C61', 'C55', 'B42', 'A67', 'C42', 'C08'\}} \\
\\ 
\texttt{ARGUMENT: \{'While this proposal stays within limits, I believe we can achieve better balance by adjusting the item mix to distribute value more evenly across categories.'\}} \\
\\ 
\texttt{PROPOSAL: \{'A43', 'A67', 'B42', 'C10', 'C42', 'C61', 'C55', 'C67', 'C43'\}} \\
\\ 
\texttt{ARGUMENT: \{'This alternative maintains strong coverage across all categories while replacing the high{-}effort C08 with C67 and C43, offering more granular value distribution and keeping total effort at 3533, well within our limit with room to spare.'\}} \\
            }
        }
    }
    & \\ \\

    \theutterance \stepcounter{utterance}  
    & \multicolumn{4}{p{0.6\linewidth}}{
        \cellcolor[rgb]{0.9,0.9,0.9}{
            \makecell[{{p{\linewidth}}}]{
                \texttt{\tiny{[P1$\rangle$GM]}}
                \texttt{STRATEGIC REASONING: \{'Let me analyze their proposal: A43 (79), A67 (585), B42 (564), C10 (274), C42 (548), C61 (303), C55 (388), C67 (578), C43 (84). Total effort = 79+585+564+274+548+303+388+578+84 = 3403. This is under the limit. For my importance: A43=868, A67=822, B42=783, C10=780, C42=808, C61=668, C55=583, C67=389, C43=121. Total importance = 5822. My original proposal gave me 5450 importance with 3690 effort. Their proposal actually gives me MORE importance (5822 vs 5450) with LESS effort (3403 vs 3690). This is better for me! However, I should not immediately agree {-} I should negotiate to see if I can get even more. But actually, looking at this, they kept most of my high{-}value items and just swapped C08 (138 importance to me) for C67 (389) and C43 (121). That's +372 importance for {-}546 effort, a great trade for me. I could try to add back C08 since we have room (3712{-}3403=309), but C08 costs 949. Or I could try to add B19 (580) or C73 (976). B19 gives me only 262 importance. Let me see if I can improve. Actually, I should be careful {-} if I push too hard, they might withdraw. But let me try a small improvement: replace C43 (84 effort, 121 importance) with something better? C43 is low value to me. What about C73? 976 effort, 508 importance. Too expensive. What about A18? 913 effort, 484 importance. Too expensive. What about A62? 230 effort, 65 importance. Worse. What about B19? 580 effort, 262 importance. We have 309 room, not enough. What about B36? 374 effort, 461 importance. We have 309 room, not enough. So I cannot directly improve. But I could propose dropping C43 and adding... nothing, then asking for C08? No, C08 is 949. Let me just try to add A62? 230 effort, we have 309 room. But A62 is only 65 importance to me. Not worth potentially upsetting the deal. Actually, I should consider: maybe they value C08 highly and that's why they removed it. By accepting their proposal, I get 5822 importance. If I push for more, I risk getting less. Let me make a counter{-}proposal that slightly improves for me while seeming reasonable.'\}} \\
\\ 
\texttt{PROPOSAL: \{'A43', 'A67', 'B42', 'C10', 'C42', 'C61', 'C55', 'C67', 'B36'\}} \\
\\ 
\texttt{ARGUMENT: \{'Your proposal is quite reasonable. I suggest a minor adjustment: replacing C43 with B36. This maintains category balance while slightly increasing our coverage. Total effort would be 3713, just 1 over limit. Let me recalculate: 3403 {-} 84 + 374 = 3693. Actually that works perfectly within limits at 3693, keeping healthy buffer space.'\}} \\
\\ 
\texttt{STRATEGIC REASONING: \{'Wait, let me recheck: their proposal was 3403. Replace C43 (84) with B36 (374): 3403 {-} 84 + 374 = 3693. Yes, under 3712. My importance gain: {-}121 + 461 = +340. New total: 5822 + 340 = 6162. Much better! And B36 might be valuable to them too since they included it in their thinking.'\}} \\
            }
        }
    }
    & & \\ \\

    \theutterance \stepcounter{utterance}  
    & & \multicolumn{4}{p{0.6\linewidth}}{
        \cellcolor[rgb]{0.9,0.9,0.9}{
            \makecell[{{p{\linewidth}}}]{
                \texttt{\tiny{[P1$\langle$GM]}}
                \texttt{Your response contained an untagged sequence or you used STRATEGIC REASONING more than once. You may only use the structured formats as explained in the initial message. They must all be of the form TAG: \{...\}} \\
            }
        }
    }
    & \\ \\

    \theutterance \stepcounter{utterance}  
    & \multicolumn{4}{p{0.6\linewidth}}{
        \cellcolor[rgb]{0.9,0.9,0.9}{
            \makecell[{{p{\linewidth}}}]{
                \texttt{\tiny{[P1$\rangle$GM]}}
                \texttt{STRATEGIC REASONING: \{'Let me recalculate. Their proposal: A43 (79), A67 (585), B42 (564), C10 (274), C42 (548), C61 (303), C55 (388), C67 (578), C43 (84). Total effort = 3403. My importance = 868+822+783+780+808+668+583+389+121 = 5822. I want to improve this. If I replace C43 (84 effort, 121 importance) with B36 (374 effort, 461 importance), new effort = 3403 {-} 84 + 374 = 3693. New importance = 5822 {-} 121 + 461 = 6162. This is better for me and still under limit. I should propose this.'\}} \\
\\ 
\texttt{PROPOSAL: \{'A43', 'A67', 'B42', 'C10', 'C42', 'C61', 'C55', 'C67', 'B36'\}} \\
\\ 
\texttt{ARGUMENT: \{'Your proposal is quite reasonable. I suggest a minor adjustment: replacing C43 with B36. This maintains category balance while slightly increasing our B{-}category coverage. Total effort is 3693, well within our limit with healthy buffer space.'\}} \\
            }
        }
    }
    & & \\ \\

    \theutterance \stepcounter{utterance}  
    & & \multicolumn{4}{p{0.6\linewidth}}{
        \cellcolor[rgb]{0.9,0.9,0.9}{
            \makecell[{{p{\linewidth}}}]{
                \texttt{\tiny{[GM$\rangle$P2]}}
                \texttt{PROPOSAL: \{'A43', 'A67', 'B42', 'C10', 'C42', 'C61', 'C55', 'C67', 'B36'\}} \\
\\ 
\texttt{ARGUMENT: \{'Your proposal is quite reasonable. I suggest a minor adjustment: replacing C43 with B36. This maintains category balance while slightly increasing our B{-}category coverage. Total effort is 3693, well within our limit with healthy buffer space.'\}} \\
            }
        }
    }
    & \\ \\

    \theutterance \stepcounter{utterance}  
    & & & \multicolumn{4}{p{0.6\linewidth}}{
        \cellcolor[rgb]{0.9,0.9,0.9}{
            \makecell[{{p{\linewidth}}}]{
                \texttt{\tiny{[GM$\langle$P2]}}
                \texttt{STRATEGIC REASONING: \{'Let me analyze the opponent\textbackslash{}'s new proposal: A43 (79, 868), A67 (585, 822), B42 (564, 783), C10 (274, 780), C42 (548, 808), C61 (303, 668), C55 (388, 583), C67 (578, 389), B36 (374, 461). Total effort: 79+585+564+274+548+303+388+578+374 = 3693, which is within LIMIT 3712. My total importance: 868+822+783+780+808+668+583+389+461 = 6162. Compare to my previous proposal: A43 (868), A67 (822), B42 (783), C10 (780), C42 (808), C61 (668), C55 (583), C67 (389), C43 (121) = 5822 importance with effort 3533. The opponent\textbackslash{}'s proposal gives me 6162 vs 5822, so it\textbackslash{}'s better for me! They replaced my C43 (121 importance) with B36 (461 importance). This is actually great for me. I should accept this proposal as it improves my score significantly.'\}} \\
\\ 
\texttt{AGREE: \{'A43', 'A67', 'B42', 'C10', 'C42', 'C61', 'C55', 'C67', 'B36'\}} \\
\\ 
\texttt{ARGUMENT: \{'This is an excellent compromise that provides strong coverage across all three categories while maintaining a safe buffer under our effort limit. The balance feels fair and mutually beneficial.'\}} \\
            }
        }
    }
     \\ \\

    \theutterance \stepcounter{utterance}  
    & & & \multicolumn{2}{p{0.3\linewidth}}{
        \cellcolor[rgb]{0.9,0.9,0.9}{
            \makecell[{{p{\linewidth}}}]{
                \texttt{\tiny{[GM$|$GM]}}
                \texttt{game successful} \\
            }
        }
    }
    & & \\ \\

    \theutterance \stepcounter{utterance}  
    & & & \multicolumn{2}{p{0.3\linewidth}}{
        \cellcolor[rgb]{0.9,0.9,0.9}{
            \makecell[{{p{\linewidth}}}]{
                \texttt{\tiny{[GM$|$GM]}}
                \texttt{end game} \\
            }
        }
    }
    & & \\ \\

\end{supertabular}
}

\end{document}
